\documentclass[adam,pdf,slideColor]{prosper}
\usepackage{graphicx}
\usepackage{pifont}
\usepackage{xspace}
\usepackage{alltt}
\def\occampi{{occam-\Pisymbol{psy}{112}}\xspace}
\def\picalculus{{\Pisymbol{psy}{112}-calculus}\xspace}
\definecolor{Red}{rgb}{0.6,0.0,0.0}
\definecolor{Green}{rgb}{0.0,0.6,0.0}
\definecolor{Blue}{rgb}{0.0,0.0,0.6}
\definecolor{Purple}{rgb}{0.6,0.0,0.6}
\definecolor{Brown}{rgb}{0.5,0.5,0.0}
\def\keyword#1{{\color{Green}#1}}
\def\operator#1{{\color{Purple}#1}}
\def\highlight#1{{\color{Brown}#1}}

\Logo(-1.0,7.3){\includegraphics[width=3cm]{kent.eps}}

\title{Exploring nanopass compilation in Haskell}
\author{Adam Sampson}
\email{{\tt ats1@kent.ac.uk}}
\institution{University of Kent\\ {\tt http://www.cs.kent.ac.uk/}}
\slideCaption{FCO}

\begin{document}

\maketitle

\begin{slide}{Why?}
\begin{itemize}
\item We've been looking at using Scheme
\begin{itemize}
\item Popular in the States, less so over here
\item Good libraries for compilation
\end{itemize}
\item Lots of Haskell users in the lab already
\begin{itemize}
\item \ldots including all our undergrads
\item Also heavily used for compiler work
\end{itemize}
\item Should we use Haskell instead?
\end{itemize}
\end{slide}

\begin{slide}{What I've been doing}
\begin{itemize}
\item Building bits of FCO: a nanopass occam compiler in Haskell
\begin{itemize}
\item Functional Compiler for occam
\end{itemize}
\item This is an exploration -- a ``spike solution''
\item Checking that all the things we want to do are possible
\item Warning: may contain traces of handwaving
\item I'll assume some knowledge of occam and Haskell; if anything's
  not clear, please shout
\end{itemize}
\end{slide}

\begin{slide}{Haskell}
\begin{itemize}
\item Mature purely-functional language
\item Open spec; several implementations
\item Powerful static type system
\item Built-in pattern matching
\item Good support for monadic programming
\item Lazy evaluation
\item Supports lightweight concurrency
\item Nice syntax (for occam programmers)
\item Lots of cool recent work (some of which I'll show later)
\item See the ``History of Haskell'' HoPL paper
\end{itemize}
\end{slide}

\begin{slide}{Existing work}
\begin{itemize}
\item occam2.1 language spec -- with a BNF syntax
\item 42 -- the model for this, in Scheme
\item occ21 -- the only full existing implementation, in C
\item JHC -- Haskell to C, in Haskell
\item Pugs -- Perl 6 to various, in Haskell
\item (Also GHC, Mincaml and a few others)
\end{itemize}
\end{slide}

\begin{slide}{A nanopass refresher}
\begin{itemize}
\item Parse language into AST
\item Do many small passes over the AST
\begin{itemize}
\item Checks
\item Annotations
\item Transformations
\end{itemize}
\item Output
\end{itemize}
\end{slide}

\begin{slide}{FCO's target}
\begin{itemize}
\item Translate occam to idiomatic C
\begin{itemize}
\item Use CIF for concurrency
\item Bugbear: calculating stack usage
\end{itemize}
\item Whole-program compilation
\begin{itemize}
\item Allows whole-program optimisations and checks
\item Can still do separate parsing/checks/early passes; just do final
  translation at ``link time''
\end{itemize}
\item FCO is not finished, but I believe the approach is sound
\end{itemize}
\end{slide}

\begin{slide}{Challenges}
\begin{itemize}
\item Parsing
\item Data structures
\item Writing transformations
\item Driver and IO
\item Extensibility
\item Debugging and tracing
\end{itemize}
\end{slide}

\begin{slide}{Parsing}
\begin{itemize}
\item 42 hasn't tackled this yet
\item occam2.1 has some oddities:
\begin{itemize}
\item Lots of lookahead needed
\item Indentation-based syntax
\item Odd line-continuation rules
\item Ambiguities
\item Left-recursive productions
\end{itemize}
\item occ21's parser keeps lots of state; we don't want to do that
\end{itemize}
\end{slide}

\begin{slide}{Parsing by passes}
\begin{itemize}
\item Break down ``parsing'' into several passes:
\begin{itemize}
\item Execute (some) preprocessor instructions
\item Detect and rejoin continuation lines
\item Convert indentation changes to tokens
\item \ldots then just use a regular parser
\item Resolve ambiguities later
\end{itemize}
\item (\verb|occamdoc| does much of this already)
\end{itemize}
\end{slide}

\begin{slide}{Parsing by Parsec}
\begin{itemize}
\item A combinator-based parsing library for Haskell
\item Productions look pretty much like BNF, e.g.:
\begin{verbatim}
specification = declaration
                <|> abbreviation
                <|> definition
sequence = do { sSEQ ; eol ; indent ;
                ps <- many1 process ;
                ...
\end{verbatim}
\item Uses Prolog-style backtracking and cuts -- no lookahead problems
\item All done using monads
\end{itemize}
\end{slide}

\begin{slide}{Monads}
\begin{itemize}
\item A data type that wraps a value
\item A monadic function can return a value and/or change the wrapper
\item For example:
\begin{itemize}
\item \verb|State| monad contains a user-defined state value
\item \verb|IO| monad has ``state of the world'' (so IO functions
  can interact with the world)
\item \verb|Parser| monad (from Parsec) keeps track of file location,
  backtracking, etc.
\end{itemize}
\end{itemize}
\end{slide}

\begin{slide}{Monad combinators}
\begin{itemize}
\item Combinators let you chain monadic functions together
\item Each monad defines how the ``bind'' operator (a standard
  combinator) works -- e.g. pass state through, backtrack upon
  failure, etc.
\item Haskell has special syntax (\verb|do| blocks) for bind
\item Monads can define their own combinators (e.g. \verb,<|>, in Parsec)
\end{itemize}
\end{slide}

\begin{slide}{Monad example}
\begin{itemize}
\item Function in \verb|State| monad that returns unique ID
\item The \verb|get| and \verb|put| functions come from the monad
\item \verb|return| wraps a value in the monad
\item \verb|do| block chains the three functions together
\end{itemize}
\begin{verbatim}
-- Define our own monad based on State.
-- The state is a single Int.
type UniqueState t = State Int t

-- This returns an Int in our monad.
getID :: UniqueState Int
getID = do  counter <- get
            put (counter + 1)
            return counter
\end{verbatim}
\end{slide}

\begin{slide}{Ambiguities}
\begin{itemize}
\item Two ambiguities in occam 2.1 syntax as specified in the
  manual
\item \verb|c ! x ; y|
\begin{itemize}
\item Is \verb|x| a tag or a variable name?
\end{itemize}
\item \verb|foo[x]|
\begin{itemize}
\item Is this an array element, or\ldots
\item \ldots is this retyping the array literal \verb|[x]| to the
  user-defined type \verb|foo|?
\item You can't retype an array -- fix the grammar to reflect this
\end{itemize}
\item Any more? What about \occampi?
\end{itemize}
\end{slide}

\begin{slide}{Data structures}
\begin{itemize}
\item Represent an AST node
\item Metadata
\begin{itemize}
\item Source position
\item Annotations from passes
\end{itemize}
\item In occ21 and 42, a record for each node type
\item In Haskell, I've tried two approaches
\end{itemize}
\end{slide}

\begin{slide}{Data structures: ``parse tree''}
\begin{itemize}
\item A catch-all \verb|Node| ADT
\item \verb|Meta| contains metadata
\begin{verbatim}
data Node = (Meta, NodeType)

data NodeType = Seq [Node]
                | AltGuard Node Node
                | Name String
                | ...
\end{verbatim}
\item Pattern matching's easy
\item Straightforward to write traversal code
\item This feels a bit non-Haskell-ish
\end{itemize}
\end{slide}

\begin{slide}{Data structures: ``AST''}
\begin{itemize}
\item One ADT per production (with simplifications)
\begin{verbatim}
data Process = Seq Meta [Process]
               | Alt Meta [AltGrd]
               | Assign [Var] ExpList
               | Skip Meta | ...

data AltGrd = AltGrd Meta Guard Process
              | ...
\end{verbatim}
\item This is what other compilers do
\item \ldots but writing code to process it is harder
\item (I'll come back to this later)
\end{itemize}
\end{slide}

\begin{slide}{Transformations}
\begin{itemize}
\item Walk over the tree, finding interesting bits of it and Doing
  Stuff to them
\item Some bits are purely functional; some must carry state
\item Checks and annotations are just transformations that don't
  change the tree
\item 42 does this using PLT's pattern-matching library
\item Can we use Haskell's built-in pattern-matching?
\end{itemize}
\end{slide}

\begin{slide}{Pattern-matching}
\begin{itemize}
\item Yes! At least, it's trivial for the \verb|Node| data structure:
\begin{verbatim}
myPass n = case n of
  Seq ps -> doSeq (map myPass ps)
  Name s -> doName s
  ...
\end{verbatim}
\item Writing the boilerplate like that for every pass is rather
  boring (and hard to extend), though
\end{itemize}
\end{slide}

\begin{slide}{Base passes}
\begin{itemize}
\item Rather than direct recursion, have a
  ``base pass'' function
\begin{verbatim}
basePass top n = case n of
  Seq ps -> Seq (map top ps)
  Name s -> Name s
\end{verbatim}
\item Then every pass can use that:
\begin{verbatim}
doStuff n = case n of
  Thing a -> doThing (pt a)
  otherwise -> pt a
  where  pt = basePass doStuff
\end{verbatim}
\item Slightly fancier glue for error handling, etc.
\end{itemize}
\end{slide}

\begin{slide}{Generated base passes}
\begin{itemize}
\item No need to write the base passes by hand
\item Generate them automatically from the data type definition (using
  Template Haskell, or simpler approaches)
\item Can generate several base passes for different subsets of the
  allowed productions
\end{itemize}
\end{slide}

\begin{slide}{The bad news}
\begin{itemize}
\item However, this all only works because we're restricted to one
  data type
\item And complicated passes get very messy\ldots
\item How do we do this for the ``AST'' data types?
\item We'd need to write \emph{generic} functions that would work on
  several data types
\item Fortunately\ldots
\end{itemize}
\end{slide}

\begin{slide}{Scrap Your Boilerplate}
\begin{itemize}
\item The ``SYB'' Generics package lets you turn:
\begin{verbatim}
foo :: Process -> Process
\end{verbatim}
into:
\begin{verbatim}
foo :: Typeable a => a -> a
\end{verbatim}
\ldots a function you can apply to any \verb|Typeable| type
\item A \emph{really} cute hack to provide introspection and dynamic
  typing (\verb|cast|) in Haskell
\end{itemize}
\end{slide}

\begin{slide}{An aside on typeclasses}
\begin{itemize}
\item Typeclasses are really nothing like OO classes
\item A typeclass is an interface: a set of functions
\begin{itemize}
\item The \verb|Show| typeclass provides the \verb|show| function
\item \verb|Eq| provides equality tests
\item \verb|Typeable| provides ``what's the type of this value?''
\end{itemize}
\item If a type is an instance of a typeclass, then it has
  implementations of all that typeclass's functions for that type
\item The Haskell compiler can ``derive'' instances of the built-in
  typeclasses (including \verb|Typeable| and \verb|Data|) automatically
\end{itemize}
\end{slide}

\begin{slide}{Scrap More Boilerplate}
\begin{itemize}
\item These functions do not automatically recurse
\item Instead, you get \verb|gmapT| -- ``generic map'', which peers
  carefully at the value's type, and maps across all the type
  constructor arguments
\item For example:
\end{itemize}
\begin{verbatim}
data Thing = Thing Foo Bar Baz
  deriving (Typeable, Data)
v :: Thing                 -- some Thing
t :: Typeable a => a -> a  -- a transform
f = gmap t v
f' (Thing x y z) = Thing (t x) (t y) (t z)
\end{verbatim}
\begin{itemize}
\item \verb|f| and \verb|f'| are equivalent -- but f is itself generic.
\end{itemize}
\end{slide}

\begin{slide}{Scrap The Whole Boiler}
\begin{itemize}
\item Using this, you can write a function that'll recurse across an
  entire data structure, applying the transform wherever it can
\item \ldots and, actually, they give it to you -- \verb|everywhere|
\item You also get functions to extend a generic transform with a
  type-specific case
\begin{verbatim}
doAny = doGen `extT` doPar
              `extT` doSeq ...
\end{verbatim}
\item There are monadic and query (i.e. returning a result rather than
  doing a transformation) versions of all these facilities too
\end{itemize}
\end{slide}

\begin{slide}{Scrapping My Boilerplate}
\begin{itemize}
\item By making \verb|Node| derive \verb|Data|, we can do away with
  the base passes in favour of generic functions
\item This works for any set of ADTs
\item For example:
\begin{verbatim}
doName :: Name -> Name
doName (Name s) =
  Name [if c == '.' then '_' else c
        | c <- s]

cIshNamesPass = everywhere (mkT doName)
\end{verbatim}
\item This can be applied to a \verb|Name|, \verb|Process|,
  \verb|Declaration|, etc. and works recursively
\end{itemize}
\end{slide}

\begin{slide}{Scrap Your Boilerstate}
\begin{itemize}
\item Passes that need state can use the \verb|State| monad
\item For example, tracking variable scope
\item I have a marvellous example of this but this slide is too small
  to contain it -- see the real code
\item The ``unique names'' pass is rather more concise in FCO than
  in 42 owing to the use of generics
\end{itemize}
\end{slide}

\begin{slide}{Driver and IO}
\begin{itemize}
\item The code that sets up and runs all the bits of the compiler, and
  interacts with the OS
\item Dead easy in Haskell -- very good OS interfacing
\item Worth cribbing from for \occampi's standard libraries\ldots
\item We have lists of passes to apply (like 42)
\item Passes are functions from one tree to another
\item The tree is pretty-printed after each pass (in debug mode)
\end{itemize}
\end{slide}

\begin{slide}{Extensibility}
\begin{itemize}
\item Want to be able to add:
\begin{itemize}
\item New AST node types
\item New passes
\end{itemize}
\item Also accessibility, really
\begin{itemize}
\item Undergrads should be able to write extensions
\end{itemize}
\end{itemize}
\end{slide}

\begin{slide}{Extensibility problems}
\begin{itemize}
\item Slotting in new passes is pretty trivial
\item Creating new types of nodes is harder
\item You can't extend an ADT
\item More generally, transformations can't replace a value with a
  different type of value
\item Several ways we could get around this\ldots
\end{itemize}
\end{slide}

\begin{slide}{Extensibility: ``plan ahead''}
\begin{itemize}
\item We could\ldots
\item Include everything you could ever possibly want in your data
  types ahead of time
\item Doable with \verb|Node|
\begin{itemize}
\item Early version of FCO extracted comment annotations in
  the \verb|Node| definition to decide which passes each node type was
  valid in
\item Metaprogramming -- you can do it (Template Haskell), but it's a
  bit experimental\ldots
\end{itemize}
\item Really nasty with lots of ADTs
\end{itemize}
\end{slide}

\begin{slide}{Extensibility: ``whatEVERRRRR''}
\begin{itemize}
\item We could\ldots
\item Use a very general data type
\begin{verbatim}
data Node = Node String Meta [Node]
\end{verbatim}
\item Emulating dynamic typing
\item So why use a statically-typed language?
\end{itemize}
\end{slide}

\begin{slide}{Extensibility: ``from time to time''}
\begin{itemize}
\item We could\ldots
\item Have several sets of ADTs
\item When we need to change the language, have a big transformation
  pass between them
\item This is what JHC/Pugs do (and what FCO does currently)
\item Not really nanopass any more
\begin{itemize}
\item The transformation passes aren't necessarily entirely mechanical
\end{itemize}
\item Awkward to extend
\end{itemize}
\end{slide}

\begin{slide}{Extensibility: ``deep magic''}
\begin{itemize}
\item We could\ldots
\item Use typeclasses and existential types (a la \verb|HList|)
\item Doable but very messy
\item Could require generating class instances automatically (using
  DrIFT or similar) as ``glue''
\item Not clear how this interacts with generics
\item Makes pattern-matching awkward (impossible?)
\end{itemize}
\end{slide}

\begin{slide}{Extensibility: the big worry}
\begin{itemize}
\item I feel like I'm fighting the type system\ldots
\item Is the nanopass approach fundamentally incompatible with static
  typing?
\item Or, alternately, is the Haskell type system insufficiently
  powerful to express it cleanly?
\begin{itemize}
\item Given the Haskell type system is one of the most powerful out
  there, these two are essentially equivalent
\end{itemize}
\item Is there some clean approach to this that I'm not seeing?
\end{itemize}
\end{slide}

\begin{slide}{Debugging}
\begin{itemize}
\item Even with static typing, many programs do not work right first
  time \verb|:-)|
\item Haskell lacks debugging facilities
\item Lazy evaluation can make debugging harder
\item At the very least, we need tracing (as 42 has) -- ``\verb|printf|
  debugging''
\item To do that, you need to use another monad
\end{itemize}
\end{slide}

\begin{slide}{Monad transformers}
\begin{itemize}
\item Much of our code will be running in a monad already
\item \ldots but if you want to use more than one monad?
\item Monad transformers let you ``stack up'' monads (e.g. State
  around IO)
\item Getting at actions in the ``inner'' monads is slightly awkward
  (\verb|liftM|)
\item \ldots and we want tracing, plus error-handling, plus the monads
  each pass needs -- perhaps four stacked monads in some code
\item Lots of wrapper functions
\begin{verbatim}
trace' = liftM . liftM . liftM . trace
\end{verbatim}
\end{itemize}
\end{slide}

\begin{slide}{Tracing}
\begin{itemize}
\item Awkward to insert tracing into monadic code
\item \ldots but it's much more awkward to trace (or support failure
  in) functional code, because you have to first make it monadic
\item Can't be hidden with wrappers
\item \ldots so, in practice, everything becomes monadic
\item A downside of the purely functional approach
\end{itemize}
\end{slide}

\begin{slide}{Conclusions}
\begin{itemize}
\item Haskell's a really elegant language\ldots
\item \ldots but I don't think it's suited to this particular problem
\item A number of useful things have come out of FCO, though\ldots
\end{itemize}
\end{slide}

\begin{slide}{Lessons}
\begin{itemize}
\item Decomposing the parser makes it simpler
\item Use a backtracking parser
\item Look at the full language
\begin{itemize}
\item Subtleties can require big changes
\item Can test on real code
\end{itemize}
\item Take advantage of generic/introspective programming (perhaps
  with Scheme record annotations)
\item Consider doing occam-to-C translation
\begin{itemize}
\item At a higher level than the VM
\end{itemize}
\end{itemize}
\end{slide}

\begin{slide}{And finally}
\begin{itemize}
\item Any questions?
\end{itemize}
\end{slide}

\begin{slide}{References}
\begin{itemize}
\item The code's on my web site: \\
\verb|http://offog.org/darcs/research/fco/|
\item Wikibooks Haskell tutorial: \\
\verb|http://en.wikibooks.org/wiki/Haskell|
\item History of Haskell: \\
{\tiny \verb|http://haskell.org/haskellwiki/History_of_Haskell|}
\item JHC: \\
\verb|http://repetae.net/john/computer/jhc/|
\item Pugs: \\
\verb|http://www.pugscode.org/|
\item Scrap Your Boilerplate: \\
\verb|http://www.cs.vu.nl/boilerplate/|
\end{itemize}
\end{slide}

\end{document}
